\documentclass[11pt]{article}

\usepackage[T1]{fontenc}
\usepackage{lmodern}
\usepackage[a4paper,margin=29mm]{geometry}
\usepackage{amsmath,amssymb,amsthm,mathtools}
\usepackage{enumitem}
\usepackage{xcolor}
\usepackage[colorlinks=true,linkcolor=blue!55!black,urlcolor=blue!55!black]{hyperref}
\usepackage[nameinlink,capitalize,noabbrev]{cleveref}

\allowdisplaybreaks
\setlength{\parindent}{0pt}
\setlength{\parskip}{0.5em}

\newtheorem{theorem}{Theorem}[section]
\newtheorem{lemma}[theorem]{Lemma}
\newtheorem{proposition}[theorem]{Proposition}
\newtheorem{conjecture}[theorem]{Conjecture}
\theoremstyle{remark}
\newtheorem{remark}[theorem]{Remark}

\newcommand{\E}{\mathbb E}
\newcommand{\Pp}{\mathbb P}
\newcommand{\one}{\mathbf 1}
\newcommand{\opt}{\operatorname{OPT}}
\newcommand{\cost}{\operatorname{cost}}
\newcommand{\divg}{\operatorname{div}}
\newcommand{\disc}{\operatorname{disc}}

\title{The $n^{1/5}$ gateway for the asymmetric clairvoyance gap}
\author{Research attack note}
\date{22 August 2026}

\begin{document}
\maketitle

\begin{abstract}
This note audits the proposed $O(n^{1/4}\sqrt{\log n})$ policy and attacks
the next exponent.  One improvement is unconditional: jointly rounding
$(2x,2y)$ removes the independent endpoint-sampling error
$\sqrt{p_*\mu}$, replacing it by deterministic $O(\log^2 n)$ public
imbalance.  This is necessary for any $n^{1/5}$ proof.

The remaining exponent is exact.  The present low-probability policy balances
a certificate term $\alpha^{-1}$ against independent-tail discrepancy
$(\alpha\mu)^{1/2}$ and therefore costs $\widetilde O(\mu^{1/3})$.
An $\widetilde O((\alpha\mu)^{1/3})$ \emph{causal fixed-inventory interval
controller} would instead give $\widetilde O(\mu^{1/4})$ and, after
probability thresholding, $\widetilde O(n^{1/5})$.  We state this gateway as
an exact conditional theorem.

We also record why four tempting substitutions do not prove the gateway:
independent tail controllers genuinely fluctuate by $\Omega(\sqrt m)$;
the existing two-arc certificate genuinely costs $\Omega(1/\alpha)$;
postponing the multiplicity decision until all activations are known has
already incurred square-root backbone cost; and grouping posterior paths
into $h$-vertex fixed-terminal packets loses precisely the factor $h$ it was
intended to save.  Thus this document does not claim an unconditional
$n^{1/5}$ theorem.  It isolates the new causal lemma that such a theorem
requires and removes an avoidable obstacle from the existing algorithm.
\end{abstract}

\section{Status and exponent ledger}

We assume $n\ge2$ and put
\[
  L:=\max\{1,\log n\}.
\]
The cases $n\le1$ are immediate.  Let $B=\opt_{\rm post}$, let
$\mu=\sum_i p_i$, and let $p_*:=\max_i p_i$ in the low-probability
subinstance.

The thinned-backbone policy in the supplied note has the bound
\begin{equation}\label{eq:old-low}
 O\!\left(
   \frac L\alpha+\sqrt{p_*\mu}
   +\sqrt{\alpha\mu L}+L
 \right)B.
\end{equation}
The polynomial part comes from
\[
  \frac L\alpha
  \quad\hbox{versus}\quad
  \sqrt{\alpha\mu L}.
\]
Optimizing $\alpha$ gives $\widetilde O(\mu^{1/3})B$.  Splitting at
probability $\tau$ then balances
\[
  \frac1\tau
  \quad\hbox{versus}\quad
  (n\tau)^{1/3},
\]
which forces the exponent $1/4$.

For a final exponent $1/5$, the low-probability policy must cost
$\widetilde O(\mu^{1/4})B$, unless the $1/\tau$ policy for the
high-probability vertices is itself improved.  The first task is to remove
$\sqrt{p_*\mu}$: at the putative threshold $\tau=n^{-1/5}$, that term alone
can be as large as $\sqrt{n\tau^2}=n^{3/10}$.

\section{A rigorous improvement: scaled joint rounding}

Fix a backbone $R$, a fractional outside-detour coupling
$x=(x_{i,ab})$, and the backbone part $y=(y_{ab})$ of the circulation.  Thus
\[
  \sum_{a,b\in R}x_{i,ab}=p_i
  \qquad\text{and}\qquad
  \Delta_{x,y}(I)=0
\]
for every outside vertex $i$ and every cyclic interval $I$ of the backbone.

\begin{lemma}[scaled joint rounding]\label{lem:scaled-rounding}
Suppose $p_i\le1/4$ for every outside vertex.  There is a randomized public
rounding, independent of the activations, which chooses one ordered pair
$(a_i,b_i)$ for every outside vertex and a nonnegative integer vector
$N=(N_{ab})$ such that every outcome satisfies
\begin{equation}\label{eq:scaled-public-balance}
 \max_I\left|
   \sum_i2p_i\chi_I(a_i,b_i)
   +\sum_{a,b\in R}N_{ab}\chi_I(a,b)
 \right|=O(\log^2 n).
\end{equation}
Moreover,
\begin{align}
 \E\!\left[
   \sum_i2p_i\bigl(d(a_i,i)+d(i,b_i)\bigr)
   +\sum_{a,b\in R}d(a,b)N_{ab}
 \right]
 =2\cost(z).                                      \label{eq:scaled-cost}
\end{align}
\end{lemma}

\begin{proof}
Run the joint randomized rounding already used for $(x,y)$, but initialize
the variables at $(2x,2y)$.  Give each outside variable the bounds
\[
  0\le \xi_{i,ab}\le2p_i,
  \qquad
  \sum_{a,b}\xi_{i,ab}=2p_i,
\]
and give each backbone variable the consecutive-integer bounds
\[
  \lfloor2y_{ab}\rfloor
  \le\eta_{ab}\le
  \lceil2y_{ab}\rceil.
\]
The initial dyadic interval equations hold because
$\Delta_{2x,2y}=0$.  Retain every outside-vertex equation and retain a
dyadic interval equation while it contains at least
$8(1+\lceil\log_2|R|\rceil)$ unfixed variables.  The same incidence count
as in the unscaled rounding leaves a nonzero null direction at every
nonterminal iteration.  Moving randomly to one of the two maximal endpoints
preserves the current vector in conditional expectation.

At termination, the outside equation and the bounds imply that exactly one
$\xi_{i,ab}$ equals $2p_i$ for every $i$.  Every $\eta$ is integral.  A
$\xi$-variable changes by at most $2p_i\le1/2$, and an $\eta$-variable
changes by at most one.  Therefore each discarded dyadic equation has final
error $O(\log n)$; decomposing a cyclic interval into $O(\log n)$ dyadic
intervals proves \eqref{eq:scaled-public-balance}.  Preservation of the
expectation of every variable and linearity of the objective prove
\eqref{eq:scaled-cost}.
\end{proof}

This replaces both the independent endpoint sampling and the subsequent
integer correction in the supplied note.  Its public mean target is exactly
$2p_i$, which matches the multiplicity controller's identity
$\E[A_iK_i]=2p_i$.

\begin{proposition}[improved proved low-probability bound]
\label{prop:improved-low}
For $p_*\le1/4$ and every $0<\alpha\le1$, the supplied thinned-backbone and
tail-controller policy, with \cref{lem:scaled-rounding} in place of
independent endpoint sampling, has expected cost
\begin{equation}\label{eq:proved-low}
 O\!\left(
   \frac L\alpha+\sqrt{\alpha\mu L}+L^2
 \right)B.
\end{equation}
Consequently its optimized bound is
\begin{equation}\label{eq:proved-low-optimized}
 O\!\left(L^2+L^{2/3}\mu^{1/3}\right)B.
\end{equation}
\end{proposition}

\begin{proof}
The thinned certificate still supplies
\[
  \cost(T_R)=O(B),\qquad
  \cost(z)=O(LB/\alpha),\qquad
  |R|=O(1+\alpha\mu)
\]
after averaging.  Fix a scaled joint-rounding outcome.  Independent
one-dimensional controllers on the head-ordered lists at the distinct tails
give expected maximum cyclic-interval error
\[
  O\!\left(\sqrt{|R|L}+L\right).
\]
The public part now contributes $O(L^2)$ deterministically, with no
$\sqrt{p_*\mu}$ term.  The active-first stack lemma therefore bounds the
expected backbone multiplicity by
\[
  O\!\left(L^2+\sqrt{\alpha\mu L}\right).
\]
The expected detour and auxiliary-chord cost is at most a constant times
$\cost(z)$ by \eqref{eq:scaled-cost} and
$\E[A_iK_i]=2p_i$.  Combining the terms proves
\eqref{eq:proved-low}.  Balancing $L/\alpha$ and
$\sqrt{\alpha\mu L}$ proves \eqref{eq:proved-low-optimized}; the cases in
which the formal optimizer is outside $(0,1]$ are absorbed by $O(L^2)$.
\end{proof}

\section{The exact \texorpdfstring{$n^{1/5}$}{n-to-the-one-fifth} gateway}

The remaining useful structure is narrower than arbitrary online vector
balancing.  The complete chord inventory and every tail list are public.
The route reveals this inventory in the active-first stack order, which may
depend on earlier activations and multiplicities, but every list is eventually
exhausted.  A chord $a\to b$ is a signed ordinary interval in backbone-prefix
coordinates.
Write $\mathcal I(R)$ for the family of cyclic intervals of $R$.

There is a sharper description.  Let $\pi_T$ concatenate the outside
vertices by cyclic tail order, using head order within each tail, and let
$\pi_H$ sort all outside vertices by cyclic head order.  For coefficients
$c=(c_i)$, write $\disc_\pi(c)$ for the maximum absolute sum over a cyclic
interval of the order $\pi$.  Since
\begin{equation}\label{eq:two-stripe-reduction}
 \sum_i c_i\chi_I(a_i,b_i)
 =\sum_{i:a_i\in I}c_i-\sum_{i:b_i\in I}c_i,
\end{equation}
and the two index sets on the right are intervals of $\pi_T$ and $\pi_H$,
respectively,
\begin{equation}\label{eq:two-stripe-bound}
 \max_{I\in\mathcal I(R)}
 \left|\sum_i c_i\chi_I(a_i,b_i)\right|
 \le \disc_{\pi_T}(c)+\disc_{\pi_H}(c).
\end{equation}
Thus it is enough to solve a causal two-stripe rounding problem.

\begin{conjecture}[causal fixed-inventory interval controller]
\label{conj:gateway}
There are constants $C,c\ge0$ with the following property.  Consider a
public fixed inventory of tail lists on a cyclically ordered set $R$ of
$m$ vertices.  Vertex $i$ has an independent activation of probability
$p_i\le1/4$, a public chord $a_i\to b_i$, and a detour length
$\ell_i\ge d(a_i,b_i)$.  If $i$ is active, its first detour copy is used
immediately and the rule chooses $K_i\in\{1,4\}$.  If $i$ is inactive, the
rule may choose $S_i\in\{0,1\}$ and create a direct surrogate chord.  Every
choice and every extra copy with tail $a$ is finalized when the list at $a$
becomes exhausted, before priority passes to auxiliary arcs or a fallback
backbone edge.  Put
\[
 C_i:=A_iK_i+(1-A_i)S_i,
 \qquad c_i:=C_i-2p_i.
\]
There is a causal rule satisfying
\begin{align}
 \E\sum_i\left(
     A_iK_i\ell_i+(1-A_i)S_i d(a_i,b_i)
   \right)
   &\le C\sum_i p_i\ell_i,
   \label{eq:gateway-cost}\\
 \E\left[\disc_{\pi_T}(c)+\disc_{\pi_H}(c)\right]
   &\le C m^{1/3}L^c.                              \label{eq:gateway-disc}
\end{align}
The rule may use the public inventory, all activations already revealed at
open tails, and fresh internal randomness, but no unrevealed activation.
\end{conjecture}

Surrogates are quantitatively affordable, although this fact alone does not
prove the conjecture.  Before probing $i$, choose any
$q_i\in[p_i,4p_i]$ from the current state.  Use one copy when $i$ is active;
when it is inactive, use a direct surrogate independently with probability
$(q_i-p_i)/(1-p_i)$.  The resulting selection has conditional mean $q_i$ and
expected movement cost at most
\[
 p_i\ell_i+(q_i-p_i)d(a_i,b_i)
 \le q_i\ell_i\le4p_i\ell_i.
\]
This gives a cost-compatible drift range $[-p_i,2p_i]$ around the public
target $2p_i$.  The missing issue is simultaneous control of both stripes
under the endogenous order, not the local budget.

The fixed-inventory and exhaustion hypotheses are essential.  A statement
for arbitrary predictable interval arrivals is false: an adaptive scheduler
can wait for a rare excursion in one coordinate and then offer only an
orthogonal singleton, freezing the excursion.  That already forces an
additive $\Omega(\log T)$ term.  More general adaptive online-interval
coloring constructions can be worse because they manufacture new interval
endpoints in response to previous choices.  Neither operation is directly
available to the fixed public inventory of the route.

\begin{theorem}[conditional $n^{1/5}$ algorithm]\label{thm:conditional}
If \cref{conj:gateway} holds, every asymmetric stochastic TSP instance has a
causal policy of expected cost
\begin{equation}\label{eq:conditional-result}
 O\!\left(n^{1/5}L^{(1+3c)/5}+L^2\right)\opt_{\rm post}.
\end{equation}
In particular, it has an $\widetilde O(n^{1/5})$ policy.
\end{theorem}

\begin{proof}
First consider a subinstance of mass $\mu$ with all probabilities at most
$1/4$.  Use the thinned certificate and scaled joint rounding.  Replace the
independent tail controllers by \cref{conj:gateway}.  By
\eqref{eq:two-stripe-bound}, the public imbalance is $O(L^2)$ and the
realized imbalance is $O(m^{1/3}L^c)$ in expectation.  Since
$m=O(1+\alpha\mu)$, the active-first execution has cost
\begin{equation}\label{eq:conditional-low-alpha}
 O\!\left(
   \frac L\alpha+(\alpha\mu)^{1/3}L^c+L^2
 \right)B.
\end{equation}
For the nontrivial parameter range, choose
\[
  \alpha=\mu^{-1/4}L^{3(1-c)/4},
\]
truncated to $(0,1]$.  This gives
\begin{equation}\label{eq:conditional-low}
 O\!\left(
   \mu^{1/4}L^{(1+3c)/4}+L^2
 \right)B.
\end{equation}

Split the original instance at
\[
  \tau:=\min\left\{\frac14,
       n^{-1/5}L^{-(1+3c)/5}\right\}.
\]
The usual fixed-tour binning policy costs $O(1/\tau)$ on the
high-probability vertices.  The low-probability subinstance has
$\mu\le n\tau$, so \eqref{eq:conditional-low} costs at most
\[
 O\!\left((n\tau)^{1/4}L^{(1+3c)/4}+L^2\right)B.
\]
The definition of $\tau$ balances its polynomial term with $1/\tau$ and
gives \eqref{eq:conditional-result}.  Concatenating the two policies at the
root and using monotonicity of the posterior optimum completes the proof.
\end{proof}

There is a logically independent gateway.  If every subinstance whose
probabilities are at least $\tau$ admitted a
$\widetilde O(\tau^{-1/2})$ policy, then balancing that policy against the
proved low bound $\widetilde O((n\tau)^{1/3})$ would also give exponent
$1/5$.  No proof of this high-probability interpolation is currently known;
the usual attempt to localize a tour into nonempty blocks fails in an
asymmetric one-way cycle, where a single nonempty block may already require a
full revolution.  One sufficient formulation would be
\[
 \opt_{\rm adapt}
 \le \operatorname{polylog}(n)
       \sqrt{\opt_{\rm prior}\opt_{\rm post}},
\]
because $\opt_{\rm prior}\le\tau^{-1}\opt_{\rm post}$ on such a subinstance.
A direct Bellman induction for this inequality encounters an unbounded
forced-first loss: in an asymmetric metric, forcing the newly active vertex
ahead of a cheap suffix can be arbitrarily more expensive than its immediate
travel distance.  That loss requires a genuinely multistep amortization.

\section{Why the immediate shortcuts fail}

\subsection{Independent tails really cost
\texorpdfstring{$\sqrt m$}{square root m}}

Take $m$ distinct tails $a_1,\ldots,a_m$, a common head $b$, and one
probability-$1/4$ outside vertex on each chord $a_j\to b$.  For a one-item
tail, marginal unbiasedness with $K\in\{1,4\}$ forces
$\Pp(K=4\mid A=1)=1/3$.  Hence the error $E=AK-1/2$ has distribution
\[
 E=\begin{cases}
   -1/2,&\text{with probability }3/4,\\
   +1/2,&\text{with probability }1/6,\\
   +7/2,&\text{with probability }1/12.
 \end{cases}
\]
It has mean zero and variance $5/4$.  If the tail controllers are
independent, the cyclic cut containing all $a_j$ but not $b$ has error
$\sum_jE_j$, and
\[
  \E\left|\sum_{j=1}^mE_j\right|=\Omega(\sqrt m).
\]
Thus neither a sharper Bernstein inequality nor a chaining argument can
turn these independent controllers into an $m^{1/3}$ construction.

\subsection{The current two-arc certificate really costs
\texorpdfstring{$1/\alpha$}{one over alpha}}

Consider the shortest-path metric of a consistently oriented cycle and a
cell whose only sampled portal is its root-side backbone vertex.  Every
allowed two-arc outside detour must leave the portal, reach its outside
vertex, and return to the same portal; its cost is one complete revolution.
A cell of a backbone sampled at rate $\alpha p_i$ has expected outside mass
$\Theta(1/\alpha)$.  The no-outside-to-outside circulation therefore pays
$\Theta(1/\alpha)$ complete revolutions although a posterior tour pays one
when the cell is nonempty.  An $\alpha^{-1/2}$ certificate cannot result from
tightening the padded-run calculation; it must retain shared
outside-to-outside chains or create persistent new portals.

\subsection{Offline multiplicity repair is too late}

Let the backbone have two vertices $a,b$.  Present $M$ independent
probability-$p$ chords in each direction, and let $X,Y$ be the two active
counts.  Suppose the probing trail uses the mandatory first copy of every
active chord but postpones all extra multiplicities.  Flow balance of this
partial trail implies that its two backbone-edge counts differ by
$X-Y+\delta$, where $|\delta|\le1$.  Consequently it has already used some
backbone edge at least $|X-Y|-1$ times, whose expectation is
$\Omega(\sqrt{Mp})$.

By contrast, with high probability $X$ and $Y$ are within constant factors.
Then the integer intervals $[X,4X]$ and $[Y,4Y]$ intersect.  An offline
Hoffman circulation assigns multiplicities in $\{1,2,3,4\}$ whose total is
the same in both directions and needs only one backbone pass.  This
separation shows that existence of the excellent offline circulation does
not imply a causal policy: the multiplicity must be selected before the
relevant fallback edge is crossed.

Inactive surrogates evade this particular counting obstruction.  After all
the activations in the example are known, $|X-Y|$ inactive chords of the
deficient orientation make the totals equal, and their expected number is
$O(Mp)$.  Selecting them only then is too late: the square-root discovery
prefix has already been traversed.  They help only if chosen while their
tails remain open, with every cyclic cut and the weighted surrogate cost
controlled simultaneously.  This is precisely the causal requirement in
\cref{conj:gateway}; an offline circulation theorem supplies no such prefix
consistency.

\subsection{Fixed-terminal packets cancel their own gain}

Let all packet vertices have probability $p$, and let a public packet contain
$s\le h$ of them.  The packet must be traversed whenever it is nonempty, an
event of probability
\[
  q_s=1-(1-p)^s.
\]
When $ph\le1$, a full packet has $q_h=\Theta(hp)$.  Packetizing a posterior
segment of length $M$ reduces its padded path count from $M$ to
$Q=\lceil M/h\rceil$.  But converting the posterior packet system into a
public cover scales a full packet's terminal flow by
$q_h/p=\Theta(h)$.  Therefore
\[
  \frac{q_h}{p}Q=\Theta(M),
\]
and the backbone repair remains $\Theta(L/\alpha)$.

Equivalently, choose a multiplicity satisfying
\[
  \E[K\mid\text{packet nonempty}]=\frac{hp}{q_s}.
\]
Then every packet contributes expected count $hp$, but the resulting bound is
\[
  h\left(1+\frac{L}{\alpha h}\right)B
  =\left(h+\frac L\alpha\right)B.
\]
The hoped-for $h$ saving disappears.  A successful chain construction must
use newly active internal vertices as persistent portals rather than collapse
the entire chain to fixed backbone terminals.

\section{Recommended next target}

The cleanest remaining proof target is \cref{conj:gateway}, not a generic
online Koml\'os theorem.  The two-stripe reduction
\eqref{eq:two-stripe-reduction} is stronger than merely observing that every
chord is an interval vector.  The reveal order is endogenous, but it is an
exhaustive adaptive linear extension of the public head-ordered tail lists.
A valid proof must use this restriction and must finalize every extra or
surrogate copy before its tail's fallback edge is committed.

Two concrete mechanisms remain plausible:
\begin{enumerate}[leftmargin=2em]
 \item a multiscale controller that coordinates several tails before any of
 their relevant fallback edges is used; or
 \item a persistent-portal execution that keeps active outside vertices in
 the virtual state and thereby retains shared outside-to-outside paths.
\end{enumerate}
Neither mechanism is supplied by the present independent-tail controller,
ordinary online signing theorems, posterior packet rounding, or the offline
Hoffman circulation.  Proving either one with the quantitative bound
$m^{1/3}\operatorname{polylog}n$ completes the $n^{1/5}$ algorithm through
\cref{thm:conditional}.

\end{document}
